\documentclass{cmspaper}
\usepackage{placeins}
\usepackage{subfig}
\usepackage{graphicx}

\usepackage[utf8]{inputenc}
\usepackage[english]{babel}
\usepackage{csquotes}
\usepackage[backend=biber,style=ieee,citestyle=numeric-comp,sorting=none]{biblatex}
%\usepackage[backend=biber,style=chem-angew,citestyle=numeric-comp,sorting=none]{biblatex}
%\ package 'overpic' for plot test manipulation
%\usepackage[backend=biber,style=alphabetic,sorting=ynt]{biblatex}
\addbibresource{jwkbiblio.bib} %Imports bibliography file

\begin{document}

%==============================================================================
% title page for few authors

\begin{titlepage}

% select one of the following and type in the proper number:
   \cmsnote{2020/000}
%  \internalnote{2005/000}
%  \conferencereport{2005/000}
   \date{1 September 2020}

  \title{Studies of ECAL Timing Reconstruction}

  \begin{Authlist}
    %A.~Author, B.~Author, C.~Author\Aref{a}
       %\Instfoot{cern}{CERN, Geneva, Switzerland}
    J.~King\Aref{a}, N.~Minafra\Aref{b}, C.~Rogan\Aref{c}
       \Instfoot{ku}{University of Kansas, Lawrence KS, USA}
       %\Instfoot{ieph}{Institute of Experimental Physics, Hepcity, Wonderland}
  \end{Authlist}

% if needed, use the following:
\collaboration{Flying Saucers Investigation Group}
%\collaboration{CMS collaboration}

  \Anotfoot{a}{At home}
  \Anotfoot{b}{At home near CERN, Geneva, Switzerland}
  \Anotfoot{b}{At home with puppies}

  \begin{abstract}
  \centering
    To be added at a late time.
    %This is a template written in LaTeX,
    %processed with {\it cmspaper.cls} document class.
    %It is based on the {\it cernart.cls} and {\it articlet.cls} classes.
    %The above information, Title and Abstract, should be readable on any 
    %computer and should be searchable on the information server. 
    %Therefore if the abstract or the title of your note contains Greek 
    %characters and mathematical symbols, please send to the CMS secretariat 
    %alternative versions which do not contain such characters.
  \end{abstract} 

% if needed, use the following:
%\conference{Presented at {\it Physics Rumours}, Coconut Island, April 1, 2005}
%\submitted{Submitted to {\it Physics Rumours}}
\note{Preliminary version}
  
\end{titlepage}

\setcounter{page}{2}%JPP

%==============================================================================
% title page for many authors
%
%\begin{titlepage}
%  \internalnote{2005/000}
%  \title{CMS Technical Note Template}
%
%  \begin{Authlist}
%    A.~Author\Iref{cern}, B.~Author\Iref{cern}, C.~Author\IAref{cern}{a},
%    D.~Author\IIref{cern}{ieph}, E.~Author\IIAref{cern}{ieph}{b},
%    F.~Author\Iref{ieph}
%  \end{Authlist}
%
%  \Instfoot{cern}{CERN, Geneva, Switzerland}
%  \Instfoot{ieph}{Institute of Experimental Physics, Hepcity, Wonderland}
%  \Anotfoot{a}{On leave from prison}
%  \Anotfoot{b}{Now at the Moon}
%
%  \begin{abstract}
%    This is a template of a CMS paper, written in LaTeX,
%    processed with {\it cmspaper.sty} style.
%    It is based on the {\it cernart.sty} and {\it articlet.sty} styles.
%    There are two versions of the title page.
%    The current one is designed for many authors.
%    The one on the previous page is for few authors.
%    Just delete the one which you do not need.
%  \end{abstract} 
%  
%\end{titlepage}
%
%==============================================================================

%\tableofcontents
\pagebreak

%\section{Outline \label{sec:outline}}
%\input{tex/outline.tex}
%\pagebreak

\section{Introduction \label{sec:intro}}
\input{tex/Intro_s1}
\FloatBarrier

\section{ECAL Time \& Energy Reconstruction \label{sec:reco}}
\input{tex/TimeEnergyReco_s2}
%\FloatBarrier
 
\subsection{ECAL Time Resolution \label{sec:reso}}
\input{tex/Method_s3}
%\FloatBarrier

\section{ECAL Timing Algorithms and OOT Pile-up Challenges \label{sec:motive}}
\input{tex/TimeAlgosMotive_s5}
%\FloatBarrier

\subsection{Ratio Time Algorithm\label{sec:algos}}
\input{tex/Ratio_s4}
%\FloatBarrier

\subsection{Monte Carlo Truth Times\label{sec:gen}}
\input{tex/GenStudies_s7}
\FloatBarrier

\section{Cross Correlation Time Algorithm\label{sec:algos}}
\input{tex/TimeAlgos_s6}
%\FloatBarrier

\pagebreak
\FloatBarrier

%\section{Summary \label{sec:conclusion}}
%\input{tex/Conclusion_s9}
%\FloatBarrier

\appendix
\section {Ratio \& CC Run 2 Time Resolution\label{sec:results}}
\input{tex/Results_s8}
\FloatBarrier

\pagebreak
\medskip
\printbibliography[heading=bibnumbered,title={References}] 
%Prints the entire bibliography with the titel "Whole bibliography"

%\clearpage
%Filters bibliography
%\printbibliography[heading=subbibintoc,type=article,title={Articles only}]
%\printbibliography[type=book,title={Books only}]
%\printbibliography%[keyword={physics},title={Physics-related only}]
%\printbibliography[keyword={latex},title={\LaTeX-related only}]

\end{document}
